\textbf{Gerald T. Dahlberg Jr.}

Independent Researcher, Sacramento, California, USA

ORCID: {[}insert ORCID iD{]}\\
Correspondence: {[}insert preferred email{]}

\textbf{Version 1.0 --- Referee-Circulation Draft}

This manuscript distinguishes derived results, numerical experiments,
and hypotheses. The full nonlinear Einstein-matter evolution remains
future work.

\section{Abstract}\label{abstract}

We formulate a coupled radial-response and toroidal-helical transport
model intended to test whether a finite-separation equilibrium can arise
dynamically rather than being imposed as a kinematic constraint. The
model is built from a covariant action containing a nonlinear radial
scalar, two transport phase fields, and a field-dependent kinetic
mixing. We define a Dynamic Equilibrium Event (DEE) as a finite, locally
stable stationary point of the radial sector under transport
backreaction. For a minimal polynomial realization, we obtain a
closed-form finite-separation branch and show that its radial stability
follows directly from the equilibrium condition. We derive the matter
stress-energy tensor and prove that its null contraction is nonnegative
for every null vector whenever the two-field kinetic matrix is positive
definite. We then study an explicit non-oscillatory transport
configuration and solve the cylindrically symmetric linearized Einstein
equations for a localized source. After a positive source rescaling of
10\^{}-6, the largest metric perturbation is of order 10\^{}-3, while
the tested co-rotating and counter-rotating null-energy contractions
remain positive. Exploratory reduced-metric calculations are also
reported to show why prescribing off-diagonal geometry without a
compatible matter source can generate energy-condition violations. The
present result is a consistent weak-field construction and a set of
nonlinear initial-data targets; it is not yet a proof of global
non-collapse in the full Einstein-matter system.

\textbf{Keywords:} general relativity; stress-energy tensor; toroidal
transport; helical transport; nonlinear stability; null energy
condition; weak-field gravity

\subsection{Classification of claims}\label{classification-of-claims}

\begin{longtable}[]{@{}
  >{\raggedright\arraybackslash}p{(\columnwidth - 4\tabcolsep) * \real{0.3333}}
  >{\raggedright\arraybackslash}p{(\columnwidth - 4\tabcolsep) * \real{0.3333}}
  >{\raggedright\arraybackslash}p{(\columnwidth - 4\tabcolsep) * \real{0.3333}}@{}}
\toprule\noalign{}
\begin{minipage}[b]{\linewidth}\raggedright
Category
\end{minipage} & \begin{minipage}[b]{\linewidth}\raggedright
Meaning in this paper
\end{minipage} & \begin{minipage}[b]{\linewidth}\raggedright
Examples
\end{minipage} \\
\midrule\noalign{}
\endhead
\bottomrule\noalign{}
\endlastfoot
Derived result & Obtained from the stated action and assumptions. &
Field equations; stress-energy tensor; finite branch; stability; NEC
theorem. \\
Numerical experiment & Computed in a reduced or weak-field model. &
Dynamic collapse toy model; parameter scans; linearized cylindrical
solve. \\
Hypothesis & Proposed behavior requiring nonlinear verification. &
Persistence of the DEE; finite curvature; global non-collapse. \\
\end{longtable}

\section{I. Introduction}\label{i.-introduction}

Gravitational collapse is ordinarily analyzed by specifying a matter
model, solving the Einstein equations, and determining whether trapped
surfaces, horizons, or singular behavior form. A recurring difficulty in
speculative non-collapse constructions is that positive separation is
imposed directly through the geometry. Such a construction can
demonstrate non-intersection by design but cannot establish that
separation emerges from dynamics.

The objective of this work is narrower and testable. We ask whether a
radial collapse degree of freedom can be coupled to directed toroidal
and helical transport so that a finite equilibrium branch appears as a
solution of the field equations of an explicit matter model. We refer to
the transition into such a state as a Dynamic Equilibrium Event. The
term \textquotesingle event\textquotesingle{} denotes a dynamical
transition, not a new causal boundary.

Two transport regimes are distinguished. The first is a nonlinear radial
response, represented by a scalar field whose effective force changes
with its local state. The second is linear non-waveform transport:
directed flow represented by constant or slowly varying field gradients
rather than by a sinusoidal carrier.
\textquotesingle Linear\textquotesingle{} therefore describes the
transport profile, not the complete Einstein-matter equations, which
remain nonlinear.

The manuscript proceeds from a covariant action, derives the matter
equations and stress tensor, solves a minimal finite-separation branch,
and proves a null-energy condition for the matter sector. An explicit
localized configuration is then used as a source for the linearized
cylindrical Einstein equations. Earlier exploratory metric ansätze are
retained only as diagnostic calculations, because they reveal how
unsupported off-diagonal geometry may demand inadmissible stress-energy.

The principal contribution is therefore not a claim of completed
nonlinear non-collapse. It is an internally consistent weak-field
framework that identifies the mathematical conditions a full nonlinear
DEE solution must satisfy.

\subsection{A. Main contributions}\label{a.-main-contributions}

\begin{itemize}
\item
  A covariant three-field action separating radial nonlinear response
  from toroidal and helical transport.
\item
  Closed-form existence and local-stability conditions for a
  finite-separation DEE branch.
\item
  A positive-definite kinetic-matrix criterion that gives an analytical
  NEC proof for every null direction.
\item
  An explicit non-waveform transport configuration with detailed angular
  energy flux and stress.
\item
  A localized, linearized cylindrical Einstein solve satisfying the
  weak-field requirement after positive rescaling.
\item
  A precise nonlinear verification program based on
  constraint-satisfying initial data, curvature invariants, and full
  null-cone testing.
\end{itemize}

\subsection{B. Scope and limitations}\label{b.-scope-and-limitations}

This work does not present a global solution of the full nonlinear
Einstein-matter equations. It does not prove avoidance of trapped
surfaces or singularities, and it does not establish experimental
realizability. The numerical values reported below are normalized
exploratory values and require independent reproduction with a
documented numerical code before journal submission. The weak-field
solution should be understood as a consistency checkpoint and
initial-data seed.

\section{II. Conventions and geometric
framework}\label{ii.-conventions-and-geometric-framework}

We use metric signature (-,+,+,+) and units c=1.
Newton\textquotesingle s constant G is retained where useful. Greek
indices run over spacetime coordinates, and hatted Latin indices denote
components in a local orthonormal frame. Repeated indices are summed.

x\^{}μ = (t, r, φ, z). (1)

The reduced cylindrical coordinates are used to develop and test the
local mechanism. A global toroidal embedding would require periodic
identification and additional regularity analysis; that extension is not
assumed in the derivations below.

The radial scalar ρ(x) represents the finite-separation or collapse
coordinate. The phase fields Θ(x) and Ψ(x) represent toroidal and
helical transport channels, respectively. The term phase does not
require periodic wave propagation; only the gradients enter the action.

\subsection{A. Dynamic Equilibrium
Event}\label{a.-dynamic-equilibrium-event}

A DEE is defined locally by a finite radial value at which the radial
field is stationary and stable under small perturbations.

ρ = ρ\_* \textgreater{} 0, ∇\_t ρ\textbar\_* = 0, ∇\_t∇\_t ρ\textbar\_*
= 0. (2)

For a homogeneous local reduction, stability requires a positive
effective radial mass squared. In the full field problem this condition
is necessary but not sufficient for global stability.

\section{III. Covariant action and field
equations}\label{iii.-covariant-action-and-field-equations}

We consider Einstein gravity coupled to one radial scalar and two
kinetically mixed transport scalars.

S = ∫ d\^{}4x √(-g) {[} R/(16πG) - 1/2 (∇ρ)\^{}2 - V(ρ) - 1/2
K\_Θ(ρ)(∇Θ)\^{}2 - 1/2 K\_Ψ(ρ)(∇Ψ)\^{}2 - Λ(ρ) ∇\_μΘ ∇\^{}μΨ {]}. (3)

The functions K\_Θ, K\_Ψ, and Λ determine the transport response and
kinetic mixing. This action is an ansatz: its consequences are derived
exactly, but its microscopic origin is not specified.

\subsection{A. Radial equation}\label{a.-radial-equation}

□ρ - V\textquotesingle(ρ) - 1/2 K\_Θ\textquotesingle(ρ)(∇Θ)\^{}2 - 1/2
K\_Ψ\textquotesingle(ρ)(∇Ψ)\^{}2 - Λ\textquotesingle(ρ)∇\_μΘ∇\^{}μΨ = 0.
(4)

The transport terms therefore modify the radial force through
derivatives of the coupling functions. No centrifugal barrier is
inserted separately.

\subsection{B. Transport equations and
currents}\label{b.-transport-equations-and-currents}

∇\_μ {[}K\_Θ(ρ)∇\^{}μΘ + Λ(ρ)∇\^{}μΨ{]} = 0, (5)

∇\_μ {[}K\_Ψ(ρ)∇\^{}μΨ + Λ(ρ)∇\^{}μΘ{]} = 0. (6)

J\_Θ\^{}μ = K\_Θ∇\^{}μΘ + Λ∇\^{}μΨ, J\_Ψ\^{}μ = K\_Ψ∇\^{}μΨ + Λ∇\^{}μΘ.
(7)

Equations (5)-(7) define conserved directed transport currents. They
permit monotonic transport profiles and do not require sinusoidal
waveform solutions.

\section{IV. Stress-energy tensor and angular
stresses}\label{iv.-stress-energy-tensor-and-angular-stresses}

Variation of the matter action with respect to the inverse metric gives

T\_μν = ∇\_μρ∇\_νρ + K\_Θ∇\_μΘ∇\_νΘ + K\_Ψ∇\_μΨ∇\_νΨ + Λ(∇\_μΘ∇\_νΨ +
∇\_μΨ∇\_νΘ) + g\_μν L\_m, (8)

L\_m = -1/2(∇ρ)\^{}2 - V - 1/2K\_Θ(∇Θ)\^{}2 - 1/2K\_Ψ(∇Ψ)\^{}2 -
Λ∇\_αΘ∇\^{}αΨ. (9)

The sign in the final term of Eq. (8) follows from the definition T\_μν
= -2 δ(√-g L\_m)/(√-g δg\^{}\{μν\}).

\subsection{A. Local angular
decomposition}\label{a.-local-angular-decomposition}

T\_âb̂ = {[}{[}ε, q\_r, q\_φ, q\_z{]}, {[}q\_r, p\_r, π\_rφ, π\_rz{]},
{[}q\_φ, π\_rφ, p\_φ, π\_φz{]}, {[}q\_z, π\_rz, π\_φz, p\_z{]}{]}. (10)

The quantities most relevant to the counter-rotating null channel are
the angular energy flux q\_φ, azimuthal pressure p\_φ, and
radial-angular shear π\_rφ.

q\_φ = T\_t̂φ̂, p\_φ = T\_φ̂φ̂, π\_rφ = T\_r̂φ̂, π\_φz = T\_φ̂ẑ. (11)

\subsection{B. Directional null-energy
condition}\label{b.-directional-null-energy-condition}

For the co-rotating and counter-rotating local null vectors k\_±\^{}â =
(1,0,±1,0), the null contraction is

N\_± = T\_âb̂ k\_±\^{}â k\_±\^{}b̂ = ε + p\_φ ± 2q\_φ. (12)

ε + p\_φ ≥ 2\textbar q\_φ\textbar{} (13)

Equation (13) is necessary and sufficient for the NEC in both purely
azimuthal null directions. It is not by itself sufficient for every null
direction.

\section{V. Minimal solvable DEE
model}\label{v.-minimal-solvable-dee-model}

To obtain an explicit finite branch, we choose the regular polynomial
functions

V(ρ) = λ\_V/4 (ρ\^{}2 - v\^{}2)\^{}2, (14)

K\_Θ(ρ) = 1 + aρ\^{}2, K\_Ψ(ρ) = 1 + bρ\^{}2, Λ(ρ) = cρ\^{}2. (15)

Here λ\_V\textgreater0, v\textgreater0, and a,b,c are dimensionful or
dimensionless according to the normalization of ρ. In a complete
dimensional treatment, a reference mass scale must be restored.

X\_Θ=(∇Θ)\^{}2, X\_Ψ=(∇Ψ)\^{}2, X\_ΘΨ=∇\_μΘ∇\^{}μΨ. (16)

\subsection{A. Finite-separation
branch}\label{a.-finite-separation-branch}

ρ\_* {[} λ\_V(ρ\_*\^{}2-v\^{}2) + aX\_Θ + bX\_Ψ + 2cX\_ΘΨ {]} = 0. (17)

ρ\_*\^{}2 = v\^{}2 - {[}aX\_Θ+bX\_Ψ+2cX\_ΘΨ{]}/λ\_V. (18)

The collapse branch is ρ\_*=0. A real positive finite branch exists when
the right-hand side of Eq. (18) is strictly positive.

aX\_Θ+bX\_Ψ+2cX\_ΘΨ \textless{} λ\_V v\^{}2. (19)

\subsection{B. Radial stability}\label{b.-radial-stability}

m\_eff\^{}2 = V\textquotesingle\textquotesingle(ρ\_*) +
1/2K\_Θ\textquotesingle\textquotesingle X\_Θ +
1/2K\_Ψ\textquotesingle\textquotesingle X\_Ψ +
Λ\textquotesingle\textquotesingle X\_ΘΨ. (20)

m\_eff\^{}2 = 2λ\_V ρ\_*\^{}2 \textgreater{} 0 for ρ\_*\textgreater0.
(21)

Thus every positive equilibrium in this minimal homogeneous reduction is
locally stable in the radial scalar direction. Coupled metric and
inhomogeneous perturbations remain to be tested.

\subsection{C. Ghost-free kinetic
sector}\label{c.-ghost-free-kinetic-sector}

K(ρ) = {[}{[}K\_Θ, Λ{]}, {[}Λ, K\_Ψ{]}{]}. (22)

K\_Θ\textgreater0, K\_Ψ\textgreater0, det K = K\_ΘK\_Ψ-Λ\^{}2
\textgreater{} 0. (23)

det K = 1+(a+b)ρ\^{}2+(ab-c\^{}2)ρ\^{}4. (24)

A simple sufficient condition for global positivity in this polynomial
family is a≥0, b≥0, and c\^{}2\textless ab.

\section{VI. Explicit transport configurations and an NEC
theorem}\label{vi.-explicit-transport-configurations-and-an-nec-theorem}

We choose non-oscillatory linear-gradient configurations in a local
cylindrical region:

Θ = ω\_Θ t + m\_Θ φ + k\_Θ z, (25)

Ψ = ω\_Ψ t + m\_Ψ φ + k\_Ψ z. (26)

Single-valued scalar phases on a globally periodic angular coordinate
would normally require integer winding numbers or an angular-valued
field interpretation. Here m\_Θ and m\_Ψ are local winding parameters.

∂\_âΘ = (ω\_Θ,0,m\_Θ/r,k\_Θ), ∂\_âΨ = (ω\_Ψ,0,m\_Ψ/r,k\_Ψ). (27)

Let k\^{}â=(1,n\_r,n\_φ,n\_z), with n\_r\^{}2+n\_φ\^{}2+n\_z\^{}2=1.
Define

U\_Θ = k\^{}â∂\_âΘ, U\_Ψ = k\^{}â∂\_âΨ. (28)

T\_âb̂ k\^{}â k\^{}b̂ = K\_Θ U\_Θ\^{}2 + K\_Ψ U\_Ψ\^{}2 + 2ΛU\_ΘU\_Ψ. (29)

T\_kk = K\_Θ(U\_Θ + ΛU\_Ψ/K\_Θ)\^{}2 + (K\_Ψ-Λ\^{}2/K\_Θ)U\_Ψ\^{}2. (30)

The metric-proportional Lagrangian term vanishes in the null
contraction. Equations (23) and (30) therefore give the following
result.

\begin{quote}
\textbf{Proposition 1.} For the matter action in Eq. (3), if the
transport kinetic matrix is positive definite, then T\_μν k\^{}μ k\^{}ν
≥ 0 for every null vector k\^{}μ. The radial scalar contributes an
additional nonnegative square.
\end{quote}

\subsection{A. Illustrative normalized
point}\label{a.-illustrative-normalized-point}

For λ\_V=4, v=1, a=0.30, b=0.20, c=0.10, X\_Θ=0.60, X\_Ψ=0.40, and
X\_ΘΨ=-0.20, Eq. (18) gives ρ\_*\^{}2=0.945 and ρ\_*≈0.972. The
corresponding kinetic coefficients are K\_Θ≈1.2835, K\_Ψ≈1.189, and
Λ≈0.0945, with det K≈1.52\textgreater0. This point is illustrative and
not a fitted physical parameter set.

\section{VII. Exploratory reduced-model
studies}\label{vii.-exploratory-reduced-model-studies}

The calculations in this section are numerical experiments. They
motivated the action-derived construction but are not used as proofs.

\subsection{A. Dynamic radial-collapse toy
model}\label{a.-dynamic-radial-collapse-toy-model}

ρ¨ = -γρ˙ - μ/(ρ+ε\_0)\^{}2 + L\^{}2/(ρ+ε\_0)\^{}3. (31)

In normalized integrations, the L=0 trajectory reached a small collapse
threshold, while positive L generated a finite-radius barrier. This
established only that rotational support can oppose collapse in a
mechanics toy model; it did not derive the support from general
relativity.

\subsection{B. Prescribed-metric
diagnostics}\label{b.-prescribed-metric-diagnostics}

Several reduced off-diagonal cylindrical metrics were tested by
computing the inferred Einstein tensor and directly sampling
T\_μνk\^{}μk\^{}ν. Those trials frequently violated the NEC,
particularly for counter-rotating azimuthal null directions near the
inner radius during the DEE switch-on. The diagnostic conclusion was
that an off-diagonal metric should not be prescribed independently of a
compatible matter source.

A parameter search localized the most negative sampled null direction
near the transition and primarily opposite the angular transport. This
observation motivated the explicit matter-sector derivation and Eq.
(13).

\section{VIII. Linearized Einstein
solution}\label{viii.-linearized-einstein-solution}

We now solve for the weak gravitational field produced by the explicit
matter model. Write

g\_μν = η\_μν + h\_μν, \textbar h\_μν\textbar{} \textless\textless{} 1.
(32)

In harmonic gauge, the trace-reversed perturbation satisfies

□ h̄\_μν = -16πG T\_μν. (33)

For a stationary, translationally invariant cylindrical approximation
with radial dependence only,

(1/r) d/dr {[} r d h̄\_μν/dr {]} = -16πG T\_μν(r). (34)

d h̄\_μν/dr = -(16πG/r) ∫\_0\^{}r s T\_μν(s) ds. (35)

A Gaussian radial envelope is applied to the explicit field-gradient
source to obtain a localized finite computational domain. This envelope
is a numerical modeling choice and is not derived from Eq. (3); a
self-consistent localized nonlinear field solution remains future work.

\subsection{A. Source normalization and weak-field
checkpoint}\label{a.-source-normalization-and-weak-field-checkpoint}

The initial normalized source generated perturbations larger than unity,
invalidating the linear approximation. Because Eqs. (33)-(35) are linear
in the source, we introduced a positive overall amplitude ε\_s.

T\_μν → ε\_s T\_μν, h\_μν → ε\_s h\_μν. (36)

ε\_s = 10\^{}-6. (37)

\begin{longtable}[]{@{}
  >{\raggedright\arraybackslash}p{(\columnwidth - 4\tabcolsep) * \real{0.3333}}
  >{\raggedright\arraybackslash}p{(\columnwidth - 4\tabcolsep) * \real{0.3333}}
  >{\raggedright\arraybackslash}p{(\columnwidth - 4\tabcolsep) * \real{0.3333}}@{}}
\toprule\noalign{}
\begin{minipage}[b]{\linewidth}\raggedright
Diagnostic
\end{minipage} & \begin{minipage}[b]{\linewidth}\raggedright
Scaled value
\end{minipage} & \begin{minipage}[b]{\linewidth}\raggedright
Status
\end{minipage} \\
\midrule\noalign{}
\endhead
\bottomrule\noalign{}
\endlastfoot
Largest \textbar h\_μν\textbar{} & 1.16 × 10\^{}-3 & Weak-field pass \\
Minimum co-rotating N\_+ & 3.71 × 10\^{}-14 & Positive \\
Minimum counter-rotating N\_- & 1.69 × 10\^{}-14 & Positive \\
Minimum sampled energy density & 2.62 × 10\^{}-14 & Positive \\
\end{longtable}

These values constitute the linearized Einstein-DEE consistency
checkpoint. The reported digits reflect the exploratory computation and
should not be interpreted as precision results.

Linearized Einstein-DEE consistency test: PASS. (38)

\subsection{B. Meaning of the pass}\label{b.-meaning-of-the-pass}

The pass means that an explicit positive-definite matter source can
generate a small cylindrical metric perturbation while maintaining
positive tested directional NEC contractions. It does not prove that the
result persists under full nonlinear evolution.

\section{IX. Nonlinear initial-value
problem}\label{ix.-nonlinear-initial-value-problem}

The weak-field result supplies a seed for a nonlinear initial-data
construction. On an initial spacelike slice,

g\_μν(t\_0,r) = η\_μν + h\_μν(r), (39)

∂\_t g\_μν(t\_0,r) = Ȧ\_DEE(t\_0) h\_μν\^{}(DEE)(r), (40)

ρ(t\_0,r)=ρ\_*+δρ(r), ∂\_tρ(t\_0,r)=δρ̇(r), (41)

with corresponding data for Θ and Ψ. These data must be adjusted or
solved so that the Hamiltonian and momentum constraints hold.

\^{}(3)R + K\^{}2 - K\_ijK\^{}ij = 16πG E, (42)

D\_j(K\^{}ij - γ\^{}ijK) = 8πG S\^{}i. (43)

\subsection{A. Nonlinear pass
criteria}\label{a.-nonlinear-pass-criteria}

ρ\_min(t) = min\_r ρ(t,r) \textgreater{} 0, (44)

N\_min(t,r) = min\_\{k null\} T\_μν k\^{}μk\^{}ν ≥ 0, (45)

R, R\_μνR\^{}μν, R\_αβγδR\^{}αβγδ remain finite, (46)

constraint residuals remain convergent under resolution refinement. (47)

Only after these conditions are demonstrated over a stated simulation
interval can the finite-separation hypothesis be promoted to a nonlinear
numerical result.

\section{X. Discussion}\label{x.-discussion}

The development above separates three questions that are often
conflated. First, can a finite radial branch exist? Equation (18)
answers yes within the model. Second, can that branch be locally stable?
Equation (21) answers yes in the homogeneous radial sector. Third, can
the supporting matter satisfy the NEC? Proposition 1 answers yes
whenever the kinetic matrix is positive definite.

Those statements do not settle the gravitational question. The Einstein
equations can generate additional instabilities, constraint violations,
trapped surfaces, or curvature growth. The weak-field calculation shows
that there is no immediate inconsistency at sufficiently small source
amplitude, but the nonlinear regime is where the central hypothesis must
be tested.

The distinction between nonlinear radial response and linear
non-waveform transport is useful because it assigns different roles to
the degrees of freedom. The radial scalar determines the local
separation branch, while the transport gradients carry angular and
helical momentum. Their coupling changes the radial effective force
without requiring an independently postulated centrifugal term.

\subsection{A. Relation to angular
stress}\label{a.-relation-to-angular-stress}

In the local orthonormal frame, the counter-rotating null direction is
controlled by ε+p\_φ-2q\_φ. A prescribed metric shift can make q\_φ too
large compared with the supporting energy and pressure. In the
action-derived model, the full null contraction is instead the quadratic
form in Eq. (29), and positive definiteness of the kinetic matrix
protects every null direction.

\subsection{B. Falsifiability}\label{b.-falsifiability}

The framework can fail in several concrete ways: the initial constraints
may have no regular solution in the desired parameter region; the finite
branch may destabilize under coupled metric perturbations; the solution
may develop trapped surfaces or divergent curvature; or the NEC may fail
after nonlinear backreaction. Any of these outcomes would falsify the
strong non-collapse hypothesis for the tested model.

\section{XI. Status of results}\label{xi.-status-of-results}

\begin{longtable}[]{@{}
  >{\raggedright\arraybackslash}p{(\columnwidth - 4\tabcolsep) * \real{0.3333}}
  >{\raggedright\arraybackslash}p{(\columnwidth - 4\tabcolsep) * \real{0.3333}}
  >{\raggedright\arraybackslash}p{(\columnwidth - 4\tabcolsep) * \real{0.3333}}@{}}
\toprule\noalign{}
\begin{minipage}[b]{\linewidth}\raggedright
Result
\end{minipage} & \begin{minipage}[b]{\linewidth}\raggedright
Classification
\end{minipage} & \begin{minipage}[b]{\linewidth}\raggedright
Status
\end{minipage} \\
\midrule\noalign{}
\endhead
\bottomrule\noalign{}
\endlastfoot
Field equations from Eq. (3) & Derived & Complete \\
Stress-energy tensor & Derived & Complete \\
Finite branch Eq. (18) & Derived & Complete within minimal model \\
Radial stability Eq. (21) & Derived & Local homogeneous result \\
NEC proposition & Derived & Complete for the stated matter action \\
Dynamic collapse and metric scans & Numerical experiment &
Exploratory \\
Linearized cylindrical solution & Numerical experiment & Weak-field
consistency pass \\
Full nonlinear finite-separation evolution & Hypothesis & Not yet
tested \\
Global absence of singularities & Hypothesis & Not established \\
\end{longtable}

\section{XII. Conclusion}\label{xii.-conclusion}

We have constructed a covariant coupled-field model in which a radial
nonlinear degree of freedom interacts with toroidal and helical
transport gradients. A minimal polynomial realization admits a
closed-form finite branch, and every positive point on that branch is
locally stable in the radial scalar sector. The matter stress tensor
satisfies the NEC for all null directions whenever its kinetic matrix is
positive definite. An explicit localized source produces a consistent
cylindrical linearized Einstein response after a positive amplitude
rescaling, with metric perturbations of order 10\^{}-3 and positive
tested co- and counter-rotating directional NEC contractions.

The present work therefore establishes a mathematically coherent
weak-field DEE framework. It does not establish global non-collapse. The
next decisive calculation is a constraint-preserving nonlinear
Einstein-matter evolution that monitors the minimum separation, full
null cone, curvature invariants, trapped-surface diagnostics, and
convergence of the constraints.

\section{Acknowledgments}\label{acknowledgments}

The author used computational and language-model assistance during
conceptual development, algebra checking, numerical prototyping, and
manuscript organization. All scientific claims, interpretations, and
responsibility for the submitted work remain with the author. The author
should replace this statement as required by the target
journal\textquotesingle s current disclosure policy.

\section{Data and code availability}\label{data-and-code-availability}

The exploratory numerical calculations described in this draft should be
consolidated into a version-controlled repository before external
circulation. A submission version should provide the complete source
code, parameter files, software versions, convergence tests, and
machine-readable output used to generate every numerical result.

\section{Conflict of interest}\label{conflict-of-interest}

The author declares no conflict of interest.

\section{Appendix A: Variation of the
action}\label{appendix-a-variation-of-the-action}

This appendix records the essential variations. For a scalar f,
δ{[}(∇f)\^{}2{]}=δg\^{}\{μν\}∇\_μf∇\_νf. The kinetic-mixing term varies
symmetrically in Θ and Ψ. Integrating by parts and discarding boundary
terms gives Eqs. (4)-(6).

δS/δρ = 0 ⇒ Eq. (4), δS/δΘ = 0 ⇒ Eq. (5), δS/δΨ = 0 ⇒ Eq. (6). (A1)

Metric variation gives Eq. (8). Before submission, the author should
include the full line-by-line derivation and verify all sign conventions
against the chosen definition of the Einstein tensor.

\section{Appendix B: Proof of the NEC
proposition}\label{appendix-b-proof-of-the-nec-proposition}

For a null vector k\^{}μ, the metric-proportional term in Eq. (8)
vanishes. Define u=(k·∇Θ, k·∇Ψ)\^{}T. The transport contribution is
u\^{}T K u. The radial scalar contributes (k·∇ρ)\^{}2. If K is positive
semidefinite, both terms are nonnegative; if K is positive definite, the
transport term vanishes only when both projected gradients vanish.

T\_μνk\^{}μk\^{}ν = (k·∇ρ)\^{}2 + u\^{}T K u ≥ 0. (B1)

\section{Appendix C: Linearized cylindrical
solver}\label{appendix-c-linearized-cylindrical-solver}

For each component, Eq. (34) was integrated through the enclosed radial
source. With a finite outer boundary R\_max and boundary condition
h̄\_μν(R\_max)=0, the formal solution is

h̄\_μν(r) = 16πG ∫\_r\^{}\{R\_max\} {[}1/s{]} {[}∫\_0\^{}s u T\_μν(u)
du{]} ds. (C1)

The numerical implementation used finite radial sampling and cumulative
quadrature. The present draft reports only summary diagnostics. A
referee version requires grid-refinement studies and independent
residual evaluation of Eq. (34).

\section{Appendix D: Symbol glossary}\label{appendix-d-symbol-glossary}

\begin{longtable}[]{@{}
  >{\raggedright\arraybackslash}p{(\columnwidth - 2\tabcolsep) * \real{0.5000}}
  >{\raggedright\arraybackslash}p{(\columnwidth - 2\tabcolsep) * \real{0.5000}}@{}}
\toprule\noalign{}
\begin{minipage}[b]{\linewidth}\raggedright
Symbol
\end{minipage} & \begin{minipage}[b]{\linewidth}\raggedright
Meaning
\end{minipage} \\
\midrule\noalign{}
\endhead
\bottomrule\noalign{}
\endlastfoot
ρ & Radial separation/collapse scalar \\
Θ & Toroidal transport phase field \\
Ψ & Helical transport phase field \\
V(ρ) & Radial potential \\
K\_Θ, K\_Ψ & Field-dependent transport kinetic coefficients \\
Λ & Toroidal-helical kinetic mixing \\
ρ\_* & Finite DEE equilibrium value \\
m\_eff\^{}2 & Local radial perturbation mass squared \\
q\_φ & Azimuthal energy flux \\
p\_φ & Azimuthal pressure \\
π\_rφ & Radial-angular shear stress \\
N\_± & Co- and counter-rotating azimuthal null contractions \\
h\_μν & Linear metric perturbation \\
ε\_s & Positive source-amplitude scaling parameter \\
\end{longtable}

\section{References}\label{references}

{[}1{]} A. Einstein, ``Die Grundlage der allgemeinen
Relativitätstheorie,'' Ann. Phys. 49, 769 (1916).

{[}2{]} R. M. Wald, General Relativity (University of Chicago Press,
Chicago, 1984).

{[}3{]} S. W. Hawking and G. F. R. Ellis, The Large Scale Structure of
Space-Time (Cambridge University Press, Cambridge, 1973).

{[}4{]} E. Poisson, A Relativist's Toolkit: The Mathematics of
Black-Hole Mechanics (Cambridge University Press, Cambridge, 2004).

{[}5{]} T. W. Baumgarte and S. L. Shapiro, Numerical Relativity: Solving
Einstein's Equations on the Computer (Cambridge University Press,
Cambridge, 2010).

{[}6{]} M. Alcubierre, Introduction to 3+1 Numerical Relativity (Oxford
University Press, Oxford, 2008).

{[}7{]} L. H. Ford, ``Quantum coherence effects and the second law of
thermodynamics,'' Proc. R. Soc. A 364, 227 (1978). {[}Energy-condition
context; verify relevance before submission.{]}

{[}8{]} P. Martín-Moruno and M. Visser, ``Classical and semi-classical
energy conditions,'' in Wormholes, Warp Drives and Energy Conditions
(Springer, 2017).

{[}9{]} G. T. Dahlberg Jr., Dynamic Equilibrium Event and related
preprints, Zenodo records {[}insert verified DOI list{]}.

{[}10{]} Supplemental numerical repository {[}to be created and archived
before submission{]}.

\section{Internal referee checklist for Version
1.1}\label{internal-referee-checklist-for-version-1.1}

\begin{itemize}
\item[$\square$]
  Verify every sign in Eqs. (3)-(9) using a single stated convention.
\item[$\square$]
  Restore physical dimensions and introduce a reference mass/length
  scale.
\item[$\square$]
  Resolve global angular-field single-valuedness and winding-number
  conditions.
\item[$\square$]
  Provide complete numerical source code and reproducible environment.
\item[$\square$]
  Perform convergence and residual tests for the linearized solver.
\item[$\square$]
  Construct constraint-satisfying nonlinear initial data.
\item[$\square$]
  Add trapped-surface diagnostics to the nonlinear evolution plan.
\item[$\square$]
  Replace placeholder ORCID, email, DOI list, and repository references.
\item[$\square$]
  Obtain an independent technical read by a general-relativity
  specialist.
\item[$\square$]
  Update journal-format and AI-assistance disclosures immediately before
  submission.
\end{itemize}
